\ifdefined\gkver\else\def\gkver{student}\fi
\documentclass[\gkver]{gaokaozhenti}
\begin{document}

\chapter{2010年山东理综}
%% paperType: 理综物理部分

\begin{enumerate}
\item
%% number: 16
%% paperName: 2010年山东理综第 16 题 4 分
%% typeId: 1
%% type: 单选题
%% chapter: 直线运动
%% point: 运动图象
%% method: 无
%% score: 4
%% degree: 650
%% duplicateId: 0
%% body:
如图甲所示，物体沿斜面由静止滑下，在水平面上滑行一段距离后停止，物体与斜面和水平面间的动摩擦因数相同，斜面与水平面平滑连接。图乙中 $v$、$a$、$f$ 和 $s$ 分别表示物体速度大小、加速度大小、摩擦力大小和路程。图乙中正确的是\xzanswer{C}

\onepicture[w=3.6cm]{figs/16.png}

\fourchoices[answer=C, ispicture=true, h=2.2cm]
{figs/16a.png}
{figs/16b.png}
{figs/16c.png}
{figs/16d.png}

%% answer: C
%% memo:
\memoanswer{
对物体进行受力分析和过程分析知，在斜面和水平面受到的合力均为恒力，两段均为匀变速运动，所以 A、B 都不对；第一段摩擦力小于第二段，所以 C 正确；路程随着时间的变化，开始也是非线性变化，所以 D 错误。
}

\item
%% number: 17
%% paperName: 2010年山东理综第 17 题 4 分
%% typeId: 1
%% type: 单选题
%% chapter: 力物体的平衡
%% point: 物体系平衡
%% method: 整体与隔离
%% score: 4
%% degree: 650
%% duplicateId: 0
%% body:
如图所示，质量分别为 $m_{1}$、$m_{2}$ 的两个物体通过轻弹簧连接，在力 $F$ 的作用下一起沿水平方向做匀速直线运动（$m_{1}$ 在地面，$m_{2}$ 在空中），力 $F$ 与水平方向成 $\theta$ 角。则 $m_{1}$ 所受支持力 $N$ 和摩擦力 $f$ 正确的是\xzanswer{AC}

\onepicture[align=right, w=5.5cm]{figs/17.png}

\fourchoices[answer=AC]
{$N=m_{1}g+m_{2}g-F\sin\theta$}
{$N=m_{1}g+m_{2}g-F\cos\theta$}
{$f=F\cos\theta$}
{$f=F\sin\theta$}

%% answer: AC
%% memo:
\memoanswer{
选整体为研究对象，在水平方向整体受摩擦力和 $F$ 在水平方向的分力，所以 C 正确，D 错误；在竖直方向受支持力 $N$、重力和 $F$ 在该方向的分力，解得 $N=m_{1}g+m_{2}g-F\sin\theta$，所以 A 正确，B 错误。
}

\item
%% number: 18
%% paperName: 2010年山东理综第 18 题 4 分
%% typeId: 2
%% type: 多选题
%% chapter: 曲线运动
%% point: 人造地球卫星
%% method: 无
%% score: 4
%% degree: 650
%% duplicateId: 0
%% body:
1970年4月24日，我国自行设计、制造的第一颗人造地球卫星“东方红一号”发射成功，开创了我国航天事业的新纪元。“东方红一号”的运行轨道为椭圆轨道，其近地点 $M$ 和远地点 $N$ 的高度分别为 $439\Ukm$ 和 $2384\Ukm$，则\xzanswer{BC}

\onepicture[align=right, w=5.0cm]{figs/18.png}

\fourchoices[answer=BC]
{卫星在 $M$ 点的势能大于 $N$ 点的势能}
{卫星在 $M$ 点的角速度大于 $N$ 点的角速度}
{卫星在 $M$ 点的加速度大于 $N$ 点的加速度}
{卫星在 $N$ 点的速度大于 $7.9\Ukms$}

%% answer: BC
%% memo:
\memoanswer{
根据 $\frac{GMm}{r^{2}}=ma=m\frac{v^{2}}{R}$，得在 $M$ 点速度大于在 $N$ 点速度，根据机械能守恒，所以卫星在 $M$ 点的势能小于 $N$ 点的势能，A 错误，C 正确；根据 $v=R\omega$，得 B 正确；D 错误。
}

\item
%% number: 19
%% paperName: 2010年山东理综第 19 题 4 分
%% typeId: 2
%% type: 多选题
%% chapter: 交流电
%% point: 变压器
%% method: 无
%% score: 4
%% degree: 650
%% duplicateId: 0
%% body:
一理想变压器原、副线圈的匝数比为 $10$:$1$，原线圈输入电压的变化规律如图甲所示，副线圈所接电路如图乙所示，$P$ 为滑动变阻器的触头。\xzanswer{AD}

\onepicture[w=8.0cm]{figs/19.png}

\fourchoices[answer=AD]
{副线圈输出电压的频率为 $50\UHz$}
{副线圈输出电压的有效值为 $31\UV$}
{$P$ 向右移动时，原、副线圈的电流比减小}
{$P$ 向右移动时，变压器的输出功率增加}

%% answer: AD
%% memo:
\memoanswer{
由图甲知该交流电的周期是 $2\times10^{-2}\Us$，交流电的频率是 $50\UHz$，变压器不改变频率，所以 A 正确；因为原线圈电压的最大值是 $310\UV$，根据变压器的匝数与电压的关系，副线圈输出电压的最大值为 $31\UV$，因此 B 错误；匝数比不变，所以原、副线圈的电流比不变，C 错误；$P$ 向右移动时，负载电阻变小，根据 $P=\frac{U_{2}^{2}}{R}$，因此变压器的输出功率增加，D 正确。
}

\item
%% number: 20
%% paperName: 2010年山东理综第 20 题 4 分
%% typeId: 2
%% type: 多选题
%% chapter: 电场
%% point: 电场线
%% method: 无
%% score: 4
%% degree: 650
%% duplicateId: 0
%% body:
某电场的电场线分布如图所示，以下说法正确的是\xzanswer{BD}

\onepicture[align=right, w=5.0cm]{figs/20.png}

\fourchoices[answer=BD]
{$c$ 点场强大于 $b$ 点场强}
{$a$ 点电势高于 $b$ 点电势}
{若将一试探电荷 $+q$ 由 $a$ 点释放，它将沿电场线运动到 $b$ 点}
{若在 $d$ 点再固定一点电荷 $-Q$，将一试探电荷 $+q$ 由 $a$ 移至 $b$ 的过程中，电势能减小}

%% answer: BD
%% memo:
\memoanswer{
根据电场线的分布和方向可判断，A 错误 B 正确；若将一试电荷 $+q$ 由 $a$ 点释放，它不一定沿电场线运动到 $b$ 点，C 错误；将一试探电荷 $+q$ 由 $a$ 移至 $b$ 的过程中，电场力做正功，电势能减小，D 正确。
}

\item
%% number: 21
%% paperName: 2010年山东理综第 21 题 4 分
%% typeId: 2
%% type: 多选题
%% chapter: 电磁感应
%% point: 切割磁感线电动势
%% method: 无
%% score: 4
%% degree: 450
%% duplicateId: 0
%% body:
如图所示，空间存在两个磁场，磁感应强度大小均为 $B$，方向相反且垂直纸面，$MN$、$PQ$ 为其边界，$OO'$ 为其对称轴。一导线折成边长为 $l$ 的正方形闭合回路 $abcd$，回路在纸面内以恒定速度 $v_{0}$ 向右运动，当运动到关于 $OO'$ 对称的位置时\xzanswer{ABD}

\onepicture[align=right, w=4.0cm]{figs/21.png}

\fourchoices[answer=ABD]
{穿过回路的磁通量为零}
{回路中感应电动势大小为 $2Blv_{0}$}
{回路中感应电流的方向为顺时针方向}
{回路中 $ab$ 边与 $cd$ 边所受安培力方向相同}

%% answer: ABD
%% memo:
\memoanswer{
正方形闭合回路运动到关于 $OO'$ 对称的位置时，进出磁感线相同，所以穿过回路的磁通量为零，A 正确；根据 $E=Blv$，由右手定则可判断回路中感应电流的方向为逆时针方向，$E=Blv_{0}+Blv_{0}=2Blv_{0}$，因此 B 正确，C 错误；由左手定则可判断，回路中 $ab$ 边与 $cd$ 边所受安培力方向均向右，所以 D 正确。
}

\item
%% number: 22
%% paperName: 2010年山东理综第 22 题 4 分
%% typeId: 2
%% type: 多选题
%% chapter: 功和能
%% point: 功能原理
%% method: 无
%% score: 4
%% degree: 250
%% duplicateId: 0
%% body:
如图所示，倾角 $\theta=30^{\circ}$ 的粗糙斜面固定在地面上，长为 $l$、质量为 $m$、粗细均匀、质量分布均匀的软绳置于斜面上，其上端与斜面顶端齐平。用细线将物块与软绳连接，物块由静止释放后向下运动，直到软绳刚好全部离开斜面（此时物块未到达地面），在此过程中\xzanswer{BD}

\onepicture[align=right, w=5.5cm]{figs/22.png}

\fourchoices[answer=BD]
{物块的机械能逐渐增加}
{软绳重力势能共减少了 $\frac{1}{4}mgl$}
{物块重力势能的减少等于软绳克服摩擦力所做的功}
{软绳重力势能的减少小于其动能的增加与克服摩擦力所做功之和}

%% answer: BD
%% memo:
\memoanswer{
选物块为研究对象，细线对物块做负功，物块机械能减小，A 错误；物块由静止释放后向下运动，到软绳刚好全部离开斜面，软绳的重心下降了 $\frac{1}{2}l$，软绳重力势能共减少了 $\frac{1}{4}mgl$，所以 B 正确；根据功和能关系，细线对软绳做的功与软绳重力势能的减少等于其动能增加与克服摩擦力所做功之和，所以 D 正确，C 错误。
}

\item
%% number: 23
%% paperName: 2010年山东理综第 23 题 6 分
%% typeId: 4
%% type: 实验题
%% chapter: 电路
%% point: 测量电阻率
%% method: 无
%% score: 6
%% degree: 450
%% duplicateId: 0
%% body:
\begin{enumerate}
\item
某同学设计了如图所示的装置来探究加速度与力的关系。弹簧秤固定在一合适的木板上，桌面的右边缘固定一支表面光滑的铅笔以代替定滑轮，细绳的两端分别与弹簧秤的挂钩和矿泉水瓶连接。在桌面上画出两条平行线 $MN$、$PQ$，并测出间距 $d$。开始时将木板置于 $MN$ 处，现缓慢向瓶中加水，直到木板刚刚开始运动为止，记下弹簧秤的示数 $F_{0}$，以此表示滑动摩擦力的大小。再将木板放回原处并按住，继续向瓶中加水后，记下弹簧秤的示数 $F_{1}$，然后释放木板，并用秒表记下木板运动到 $PQ$ 处的时间 $t$。

\onepicture[align=right, w=6.5cm]{figs/23a.png}

\begin{steps}
\item
木板的加速度可以用 $d$、$t$ 表示为 $a=$\tkanswer[2cm]{$\frac{2d}{t^{2}}$}；为了减小测量加速度的偶然误差可以采用的方法是（一种即可）\tkanswer[3cm]{保持 $F_{1}$ 不变，重复实验多次测量求平均值}。
\item
改变瓶中水的质量重复实验，确定加速度 $a$ 与弹簧秤示数 $F_{1}$ 的关系。下列图像能表示该同学实验结果的是\tkanswer[1cm]{C}。

\fourchoices[ispicture=true, h=1.5cm]
{figs/23ba.png}
{figs/23bb.png}
{figs/23bc.png}
{figs/23bd.png}

\item
用加水的方法改变拉力的大小与挂钩码的方法相比，它的优点是\tkanswer[1.5cm]{bc}。

a．可以改变滑动摩擦力的大小

b．可以更方便地获取多组实验数据

c．可以比较精确地测出摩擦力的大小

d．可以获得更大的加速度以提高实验精度
\end{steps}
\item
在测定金属电阻率的实验中，某同学连接电路如图所示。闭合电键后，发现电路有故障（已知电源、电表和导线均完好，电源电动势为 $E$）：

\onepicture[align=right, w=6.5cm]{figs/23.png}

\begin{steps}
\item
若电流表示数为零、电压表示数为 $E$，则发生故障的是\tkanswer[2.5cm]{待测金属丝}（填“待测金属丝”、“滑动变阻器”或“电键”）。
\item
若电流表、电压表示数均为零，该同学利用多用电表检查故障。先将选择开关旋至\tkanswer[2.5cm]{直流电压 $10\UV$}档（填“欧姆 $\times100$”、“直流电压 $10\UV$”或“直流电流 $2.5\UmA$”），再将\tkanswer[1cm]{红}（填“红”或“黑”）表笔固定在 $a$ 接线柱，把另一支表笔依次接 b、c、d 接线柱。若只有滑动变阻器断路，则多用电表的示数依次是\tkanswer[1cm]{0}、\tkanswer[1cm]{$E$}、\tkanswer[1cm]{$E$}。
\end{steps}
\end{enumerate}
%% answer: 
\jdanswer{
\begin{enumerate}
\item
（1）$\frac{2d}{t^{2}}$；保持 $F_{1}$ 不变，重复实验多次测量求平均值；（2）C；（3）bc

\item
（1）待测金属丝；（2）直流电压 $10\UV$，红，$0$、$E$、$E$

\end{enumerate}
}
%% memo:
\memoanswer{
（1）①由于 $x=\frac{1}{2}at^{2}$，且 $x=d$，所以 $a=\frac{2d}{t^{2}}$；②由于加速度 $a=\frac{F_{1}-F_{0}}{m}$，所以 $a-F_{1}$ 图像 C 正确；③BC。

（2）①电流表示数为零，说明电路断路，又电压表示数为 $E$，说明电压表跨接的待测金属丝断路，等于将电压表直接接在电源上，所以发生故障的是待测金属丝断路。由于电路断路，用多用电表检查故障应该使用电压档，所以选择开关旋至直流电压 $10\UV$ 档。$a$ 为电源的正极，所以红表笔固定在 $a$ 接线柱。若只有滑动变阻器断路，黑表笔依次接 b、c、d 接线柱，$U_{ab}=0$，$U_{ac}=U_{ad}=E$。
}

\item
%% number: 24
%% paperName: 2010年山东理综第 24 题 15 分
%% typeId: 6
%% type: 计算题
%% chapter: 运动和力
%% point: 牛顿定律的应用
%% method: 无
%% score: 15
%% degree: 450
%% duplicateId: 0
%% body:
如图所示，四分之一圆轨道 $OA$ 与水平轨道 $AB$ 相切，它们与另一水平轨道 $CD$ 在同一竖直面内，圆轨道 $OA$ 的半径 $R=0.45\Um$，水平轨道 $AB$ 长 $s_{1}=3\Um$，$OA$ 与 $AB$ 均光滑。一滑块从 $O$ 点由静止释放，当滑块经过 $A$ 点时，静止在 $CD$ 上的小车在 $F=1.6\UN$ 的水平恒力作用下启动，运动一段时间后撤去 $F$。当小车在 $CD$ 上运动了 $s_{2}=3.28\Um$ 时速度 $v=2.4\Ums$，此时滑块恰好落入小车中。已知小车质量 $M=0.2\Ukg$，与 $CD$ 间的动摩擦因数 $\mu=0.4$。（取 $g=10\Umsq$）求：
\begin{enumerate}
\item
恒力 $F$ 的作用时间 $t$；
\item
$AB$ 与 $CD$ 的高度差 $h$。
\end{enumerate}
\onepicture[align=right, w=8.0cm]{figs/24.png}
%% answer: 
\jdanswer{
\begin{enumerate}
\item
$1\Us$

\item
$0.8\Um$

\end{enumerate}
}
%% memo:
\memoanswer{
（1）设小车在恒力 $F$ 作用下的位移为 $l$，由动能定理得 $Fl-\mu Mgs_{2}=\frac{1}{2}Mv^{2}$，由牛顿第二定律得 $F-\mu Mg=Ma$，由运动学公式得 $l=\frac{1}{2}at^{2}$，联立以上三式，代入数据得 $a=4\Umsq$，$t=\sqrt{\frac{2l}{a}}=1\Us$。

（2）滑块由 $O$ 滑至 $A$ 的过程中机械能守恒，即 $mgR=\frac{1}{2}mv_{A}^{2}$，$AB$ 段运动时间为 $t=\frac{s_{1}}{v_{A}}=\frac{s_{1}}{\sqrt{2gR}}=1\Us$，故滑块离开 $B$ 后平抛时间与小车撤掉恒力 $F$ 后运动时间相同。由牛顿第二定律得 $\mu Mg=Ma'$，由运动学公式得 $v=at-a't'$，由平抛规律得 $h=\frac{1}{2}gt'^{2}$，代入数据得 $h=0.8\Um$。
}

\item
%% number: 25
%% paperName: 2010年山东理综第 25 题 18 分
%% typeId: 6
%% type: 计算题
%% chapter: 磁场
%% point: 带电粒子在磁场中的运动
%% method: 无
%% score: 18
%% degree: 450
%% duplicateId: 0
%% body:
如图所示，以两虚线为边界，中间存在平行纸面且与边界垂直的水平电场，宽度为 $d$，两侧为相同的匀强磁场，方向垂直纸面向里。一质量为 $m$、带电量 $+q$、重力不计的带电粒子，以初速度 $v_{1}$ 垂直边界射入磁场做匀速圆周运动，后进入电场做匀加速运动，然后第二次进入磁场中运动，此后粒子在电场和磁场中交替运动。已知粒子第二次在磁场中运动的半径是第一次的二倍，第三次是第一次的三倍，以此类推。求
\begin{enumerate}
\item
粒子第一次经过电场的过程中电场力所做的功 $W_{1}$。
\item
粒子第 $n$ 次进入电场时电场强度的大小 $E_{n}$。
\item
粒子第 $n$ 次经过电场所用的时间 $t_{n}$。
\item
假设粒子在磁场中运动时，电场区域场强为零。请画出从粒子第一次射入磁场至第三次离开电场的过程中，电场强度随时间变化的关系图线（不要求写出推导过程，不要求标明坐标刻度值）。
\end{enumerate}
\onepicture[align=right, w=6.0cm]{figs/25.png}
%% answer: 
\jdanswer{
\begin{enumerate}
\item
$\frac{3mv_{1}^{2}}{2}$

\item
$\frac{(2n+1)mv_{1}^{2}}{2qd}$

\item
$\frac{2d}{(2n+1)v_{1}}$

\item
如图所示

\end{enumerate}
}
%% memo:
\memoanswer{
带电粒子在磁场中做匀速圆周运动，由 $qvB=m\frac{v^{2}}{r}$ 得 $r=\frac{mv}{qB}$，则 $v_{1}:v_{2}:\cdots:v_{n}=r_{1}:r_{2}:\cdots:r_{n}=1:2:\cdots:n$。

（1）第一次过电场，由动能定理得 $W_{1}=\frac{1}{2}mv_{2}^{2}-\frac{1}{2}mv_{1}^{2}=\frac{3}{2}mv_{1}^{2}$。

（2）第 $n$ 次经过电场时，由动能定理得 $qE_{n}d=\frac{1}{2}mv_{n+1}^{2}-\frac{1}{2}mv_{n}^{2}$，解得 $E_{n}=\frac{(2n+1)mv_{1}^{2}}{2qd}$。

（3）第 $n$ 次经过电场时的平均速度 $\overline{v_{n}}=\frac{v_{n}+v_{n+1}}{2}=\frac{2n+1}{2}v_{1}$，则时间为 $t_{n}=\frac{d}{\overline{v_{n}}}=\frac{2d}{(2n+1)v_{1}}$。

（4）如图所示。

\onepicture[w=6.5cm]{figs/25_an.png}
}

\item
%% number: 36
%% paperName: 2010年山东理综第 36 题 8 分
%% typeId: 6
%% type: 计算题
%% chapter: 气体性质
%% point: 热力学第一定律
%% method: 无
%% score: 8
%% degree: 450
%% duplicateId: 0
%% body:
一太阳能空气集热器，底面及侧面为隔热材料，顶面为透明玻璃板，集热器容积为 $V_{0}$，开始时内部封闭气体的压强为 $p_{0}$。经过太阳暴晒，气体温度由 $T_{0}=300\UK$ 升至 $T_{1}=350\UK$。
\begin{enumerate}
\item
求此时气体的压强。
\item
保持 $T_{1}=350\UK$ 不变，缓慢抽出部分气体，使气体压强再变回到 $p_{0}$。求集热器内剩余气体的质量与原来总质量的比值。判断在抽气过程中剩余气体是吸热还是放热，并简述原因。
\end{enumerate}
\onepicture[align=right, w=6.5cm]{figs/36.png}
%% answer: 
\jdanswer{
\begin{enumerate}
\item
$\frac{7}{6}p_{0}$

\item
$\frac{6}{7}$，抽气过程气体体积变大，对外做功，而温度不变内能不变，由热力学第一定律知气体应吸热。

\end{enumerate}
}
%% memo:
\memoanswer{
（1）设升温后气体的压强为 $p_{1}$，由查理定律得 $\frac{p_{0}}{T_{0}}=\frac{p_{1}}{T_{1}}$，解得 $p_{1}=\frac{7}{6}p_{0}$。

（2）抽气过程可等效为等温膨胀过程，设膨胀后气体的总体积为 $V$，由玻意耳定律得 $p_{1}V_{0}=p_{0}V$，解得 $V=\frac{7}{6}V_{0}$，设剩余气体的质量与原来总质量的比值为 $k$，由题意得 $k=\frac{V_{0}}{V}=\frac{6}{7}$，吸热。因为抽气过程中剩余气体温度不变，故内能不变，而剩余气体膨胀对外做功，所以根据热力学第一定律可知剩余气体要吸热。
}

\item
%% number: 37
%% paperName: 2010年山东理综第 37 题 4 分
%% typeId: 6
%% type: 计算题
%% chapter: 光学
%% point: 全反射
%% method: 无
%% score: 4
%% degree: 650
%% duplicateId: 0
%% body:
\begin{enumerate}
\item
渔船常利用超声波来探测远处鱼群的方位。已知某超声波频率为 $1.0\times10^{5}\UHz$，某时刻该超声波在水中传播的波动图象如图所示。

\onepicture[w=6.0cm]{\includesvg{figs/37a}}

\begin{steps}
\item
从该时刻开始计时，画出 $x=7.5\times10^{-3}\Um$ 处质点做简谐运动的振动图像（至少一个周期）。
\item
现测得超声波信号从渔船到鱼群往返一次所用时间为 $4\Us$，求鱼群与渔船间的距离（忽略船和鱼群的运动）。
\end{steps}
\item
如图所示，一段横截面为正方形的玻璃棒，中间部分弯成四分之一圆弧形状，一细束单色光由 $MN$ 端面的中点垂直射入，恰好能在弧面 $EF$ 上发生全反射，然后垂直 $PQ$ 端面射出。

\onepicture[align=right, w=3.5cm]{figs/37b.png}

\begin{steps}
\item
求该玻璃棒的折射率。
\item
若将入射光向 $N$ 端平移，当第一次射到弧面 $EF$ 上时\tkanswer[2cm]{能}（填“能”、“不能”或“无法确定能否”）发生全反射。
\end{steps}
\end{enumerate}
%% answer: 
\jdanswer{
\begin{enumerate}
\item
（1）如图所示；（2）$3000\Um$

\item
（1）$\sqrt{2}$；（2）能

\end{enumerate}
}
%% memo:
\memoanswer{
（1）①由于 $x=7.5\times10^{-3}\Um$ 处质点在负的最大位移处，其周期 $T=1\times10^{-5}\Us$，质点做简谐运动的振动图像如图所示。

\onepicture[w=6.0cm]{\includesvg{figs/37a_an}}

②由波形图可以读出波长 $\lambda=15\times10^{-3}\Um$，由波速公式 $v=\lambda f=1500\Ums$，鱼群与渔船的距离 $x=\frac{vt}{2}=3000\Um$。

（2）①发生全反射的临界角 $C=\arcsin\frac{1}{n}=45^{\circ}$，得 $n=\sqrt{2}$。②由于入射角不变，在弧面 $EF$ 上仍发生全反射。
}

\item
%% number: 38
%% paperName: 2010年山东理综第 38 题 4 分
%% typeId: 6
%% type: 计算题
%% chapter: 运动和力
%% point: 动量守恒定律
%% method: 无
%% score: 4
%% degree: 450
%% duplicateId: 0
%% body:
\begin{enumerate}
\item
大量氢原子处于不同能量激发态，发生跃迁时放出三种不同能量的光子，其能量值分别是：$1.89\UeV$，$10.2\UeV$，$12.09\UeV$。跃迁发生前这些原子分布在\tkanswer[1.5cm]{2}个激发态能级上，其中最高能级的能量值是\tkanswer[2cm]{$-1.51$}$\UeV$（基态能量为 $-13.6\UeV$）。
\item
如图所示，滑块 $A$、$C$ 质量均为 $m$，滑块 $B$ 质量为 $\frac{3}{2}m$。开始时 $A$、$B$ 分别以 $v_{1}$、$v_{2}$ 的速度沿光滑水平轨道向固定在右侧的挡板运动，现将 $C$ 无初速地放在 $A$ 上，并与 $A$ 粘合不再分开，此时 $A$ 与 $B$ 相距较近，$B$ 与挡板相距足够远。若 $B$ 与挡板碰撞将以原速率反弹，$A$ 与 $B$ 碰撞将粘合在一起。为使 $B$ 能与挡板碰撞两次，$v_{1}$、$v_{2}$ 应满足什么关系？

\onepicture[align=right, w=8.0cm]{figs/38.png}
\end{enumerate}
%% answer: 
\jdanswer{
\begin{enumerate}
\item
$2$；$-1.51$

\item
$1.5v_{2}<v_{1}\leqslant 2v_{2}$

\end{enumerate}
}
%% memo:
\memoanswer{
（1）由于氢原子的基态能量为 $-13.6\UeV$，跃迁时放出三种不同能量的光子，跃迁发生前这些原子分布在 $2$ 个激发态能级上。最高能级的能量值为 $12.09-13.6=-1.51\UeV$。

（2）设向右为正方向，$A$ 与 $C$ 粘合在一起的共同速度为 $v'$，由动量守恒定律得 $mv_{1}=2mv'$，为保证 $B$ 碰挡板前 $A$ 未能追上 $B$，应满足 $v'\leqslant v_{2}$，设 $A$ 与 $B$ 碰后的共同速度为 $v''$，由动量守恒定律得 $2mv'-\frac{3}{2}mv_{2}=\frac{7}{2}mv''$，为使 $B$ 能与挡板再次碰撞应满足 $v''>0$，得 $1.5v_{2}<v_{1}\leqslant 2v_{2}$。
}

\end{enumerate}

\end{document}
